\section{Body-bar Tay theorem}
\label{sec:body-bar}

This chapter is the forward-mode authoritative dep-graph and lemma
index for Phase~12 of the formalization. It is currently a planning
skeleton; the section headers below correspond to the Layer plan
in \texttt{notes/Phase12.md}, and the dependency-graph nodes will
be added (red) and turned green as each Layer's Lean lands.

The phase target is \textbf{Tay's theorem}~\cite{tay1984}: for
$n \ge 2$ and $d = \binom{n+1}{2}$, a multigraph $G$ on $b$ bodies
is the underlying graph of an infinitesimally rigid independent
body-bar framework in $\mathbb{R}^n$ if and only if
$|E| = d(b - 1)$ and every $m$-body sub-multigraph has at most
$d(m - 1)$ edges; equivalently, by the
Tutte--Nash-Williams~\cite{tutte1961,nashWilliams1961} tree-packing
theorem, $G$ is the edge-disjoint union of $d$ spanning trees.

The proof route follows Whiteley's matroid-union
framing~\cite{whiteley1988} (a streamlining of Tay's original
induction-based proof~\cite{tay1984}, surveyed in modern form
in~\cite{whiteley1996}). Three ingredients:
\begin{enumerate}
  \item The generic $k$-frame matroid on a multigraph is the
        $k$-fold union of the cycle matroid (Whiteley's
        Theorem~1).
  \item Edmonds' matroid-partition theorem~\cite{edmonds1965}
        characterizes independence in this union by the count
        $|E'| \le k|V'| - k$ on every non-empty subset.
        Tutte--Nash-Williams~\cite{tutte1961,nashWilliams1961}
        follows as a corollary.
  \item Body-bar realizations are obtained by specializing
        two-extensor row coefficients to standard-basis
        Pl\"ucker coordinates (Whiteley's witness
        construction).
\end{enumerate}

Whiteley's full irreducible-variety lift (Proposition~6
in~\cite{whiteley1988}, lifting independence from the
indeterminate setting to ``almost all'' body-bar realizations) is
\emph{deferred} out of Phase~12. The phase ships the
existence-of-realization form of Tay's theorem only: there exists
a body-bar realization in $\mathbb{R}^n$ that is infinitesimally
rigid (resp. independent) iff the count holds, witnessed by the
standard-basis specialization.

\subsection{Matroid-union machinery}
\label{sec:body-bar-matroid-union}

% Layer 1: vendored from apnelson1/Matroid's WIP/Union.lean +
% WIP/Submodular.lean under CombinatorialRigidity/Matroid/.
% Re-states the relevant API as cited facts: Matroid.Union of an
% indexed family, the union-indep iff existence-of-decomposition,
% the matroid-partition theorem (Edmonds), and the polymatroid
% rank formula underlying it.

\subsection{The $k$-frame matroid as a union of cycle matroids}
\label{sec:body-bar-k-frame}

% Layer 2: Whiteley's Theorem 1. Definition of the generic
% k-frame matrix F(G, X) on a multigraph; statement and proof
% that the k-frame matroid equals the k-fold union of the cycle
% matroid via the column-reorder argument.

\subsection{Tree-packing as a corollary}
\label{sec:body-bar-tree-packing}

% Layer 3: Whiteley's Corollary 3. The matroid-partition theorem
% specialized to k copies of the cycle matroid yields the count
% characterization of k-forest decomposability, recovering
% Tutte--Nash-Williams.

\subsection{Body-bar frameworks and the rigidity matrix}
\label{sec:body-bar-framework}

% Layer 4: definitions following Whiteley 1988 Section 3 and Tay
% 1984 Section 5. Body-bar framework, rigidity map, infinitesimal
% rigidity, basic API mirroring \cref{sec:frameworks} for the
% body-bar setting.

\subsection{Tay's theorem (existence-of-realization form)}
\label{sec:body-bar-tay}

% Layer 5a: Whiteley's Theorem 8 in its witness form. There
% exists an independent (resp. rigid + independent) body-bar
% framework on G in R^n iff |E'| <= d|V'| - d for all non-empty
% subsets E' (equivalently, G is the union of d edge-disjoint
% forests / spanning trees). Proof: specialize each bar's
% two-extensor coordinate to a standard-basis Plücker vector
% indexed by the bar's forest in a chosen Tutte--Nash-Williams
% partition; verify independence + rigidity directly.
